% Generated draft: requires mathematical and source review.
\documentclass[11pt]{article}
\usepackage[margin=25mm]{geometry}
\usepackage{fontspec}
\setmainfont{Latin Modern Roman}
\usepackage{amsmath,amssymb,mathtools,unicode-math}
\setmathfont{Latin Modern Math}
\usepackage{enumitem,needspace,xurl,hyperref}
\setlength{\parindent}{0pt}
\setlength{\parskip}{5pt}
\setlength{\emergencystretch}{3em}
\newcommand{\N}{\mathbb{N}}
\newcommand{\Z}{\mathbb{Z}}
\newcommand{\Q}{\mathbb{Q}}
\newcommand{\R}{\mathbb{R}}
\newcommand{\C}{\mathbb{C}}
\newcommand{\F}{\mathbb{F}}
\newcommand{\A}{\mathbb{A}}
\newcommand{\PP}{\mathbb{P}}
\newcommand{\E}{\mathbb{E}}
\newcommand{\eps}{\varepsilon}
\newcommand{\dd}{\,\mathrm{d}}
\begin{document}
\section*{UnsolvedMath ID: 700022}
\textbf{Source problem code:} catalogue entry
\subsection{Definitions and mathematical statement}
Work in the first-order language of groups \(\mathcal L_{\mathrm{grp}}=\{\cdot,{}^{-1},1\}\), with equality and without parameters. A free group on a set \(S\) is a group \(F(S)\), together with a map \(\iota:S\to F(S)\), such that every map from \(S\) to a group \(H\) extends uniquely to a homomorphism \(F(S)\to H\). Its rank is \(|S|\).

A universal sentence is a sentence of the form \(\forall x_1\cdots\forall x_n\,\varphi(x_1,\ldots,x_n)\), where \(n\geq 0\) and \(\varphi\) is a quantifier-free Boolean combination of equations between group terms. For a group \(G\), define
\[
\operatorname{Th}_{\forall}(G)=\{\sigma:\sigma\text{ is a universal }\mathcal L_{\mathrm{grp}}\text{-sentence and }G\models\sigma\}.
\]
For a nonempty class \(\mathcal F\) of groups, define its common universal theory by
\[
\operatorname{Th}_{\forall}(\mathcal F)=\{\sigma:\sigma\text{ is a universal }\mathcal L_{\mathrm{grp}}\text{-sentence and }(\forall G\in\mathcal F)\ G\models\sigma\}.
\]

\textbf{Source request (M. I. Kargapolov, 1965).} Describe the universal theory of free groups.

\textbf{Unresolved scope.} The supplied statement does not specify ranks, whether only nonabelian free groups are intended, or whether the target is the theory of an individual free group or the common theory of a class of free groups. Thus the notation above defines the possible targets without selecting one. In particular, inclusion of ranks \(0\) and \(1\), and any restriction to finite rank, remain unspecified. The word ``describe'' also does not prescribe a particular form of characterization, axiomatization, or algorithmic result. These qualifications require clarification before a unique formal target can be fixed.
\subsection{Short English statement}
Kargapolov asks for a description of the universally quantified group statements true in free groups. The source leaves the intended ranks, the individual-versus-common interpretation, and the required form of description unspecified.
\subsection{Sources}
\begin{itemize}
\item Local UnsolvedMath registry: supplied complete source record, permanent ID \(700022\), exact source problem code \texttt{catalogue entry}, titled ``Kourovka Notebook Problem 1.27.'' The record identifies M. I. Kargapolov as proposer and gives the proposal year as 1965.
\item Supplied bibliographic attribution: \emph{The Kourovka Notebook (Unsolved Problems in Group Theory)}, original problem 1.27, first issue (1965); cited edition: 21st issue (2026), arXiv identifier \texttt{1401.0300v48}. The supplied collection description names E. I. Khukhro and V. D. Mazurov as editors. This reference was supplied with the record and has not been independently consulted.
\item The registry labels the problem ``open'' and records creation and update timestamps of \texttt{2026-10-16T14:00:00Z}. These are supplied registry metadata, not independent verification of mathematical status or literature coverage.
\end{itemize}
\end{document}
